% Adapted from Dynamic Programming by Pascal Michaillat. CC BY 4.0.
% See cover.tex for sources and adaptation details.
\documentclass{article}
\usepackage{hyperref}
\begin{document}
\newcounter{tpequation}
\setcounter{section}{0}
\section{Simple deterministic problem}\label{sec:deterministic}\wcSpan*[id=ref-sec-deterministic]{}

This section introduces the key concepts of dynamic programming in a deterministic problem (a problem without randomness). Section~\href{stochastic.tex#ref-sec-stochastic}{2} shows how the key concepts apply to a stochastic problem. Section~\href{theory.tex#ref-sec-theory}{4} proposes a more formal treatment of dynamic programming.

\subsection{Consumption-saving problem}\label{subsec:PROBLEM}\wcSpan*[id=ref-subsec-PROBLEM]{}

You start life with wealth $a_0>0$. Each period $t=0,1,\dots,\infty$, you consume a quantity $c_t\ge 0$ of your wealth, which provides utility $u(c_t)$. You choose consumption to maximize your lifetime utility

\[
\sum_{t=0}^{\infty}\beta^t u(c_t),
\]
where $\beta\in[0,1)$ is the discount factor. Assume that $\lim_{c\to 0} u'(c)=\infty$, which implies that you choose to consume at least some of your wealth ($c_{t}>0$) each period. You save the amount of wealth that you do not consume. The wealth at the beginning of period $t$ is $a_t$. Wealth is invested at a constant interest rate $r$, paid at the beginning of period $t$; hence, wealth evolves according to the law of motion

\[
a_{t+1}=(1+r) a_t - c_t.
\]

\subsection{Direct approach}

Take wealth $a_0>0$ at $t=0$ as given. The direct approach is to choose two sequences of variables $\left\lbrace c_t\right\rbrace _{t=0}^{\infty}$ and $\left\lbrace a_t\right\rbrace _{t=0}^{\infty}$ to maximize

\[
\sum_{t=0}^{\infty}\beta^t u(c_t)
\]
subject to the constraints that for all $t\ge 0$,
\wcSpan*[id=equation-1]{}
\setcounter{tpequation}{0}\refstepcounter{tpequation}\label{e:constraint1}
\setcounter{tpequation}{1}\refstepcounter{tpequation}\label{e:constraint2}
\[
\begin{align*}
a_{t+1}&=(1+r) a_{t}-c_{t}\tag{1}\\
a_{t+1}&\ge 0.\tag{2}
\end{align*}
\]
To find the two sequences $\left\lbrace c_t\right\rbrace _{t=0}^{\infty}$ and $\left\lbrace a_t\right\rbrace _{t=0}^{\infty}$, we write down the Lagrangian associated with the problem:

\[
\mathcal{L}=\sum_{t=0}^{\infty}\beta^t \left\lbrace u(c_t)-\lambda_{t} \left[a_{t+1}-(1+r) a_{t}+c_{t}\right]-\mu_{t} a_{t+1}\right\rbrace ,
\]
where $\left\lbrace \lambda_t\right\rbrace _{t=0}^{\infty}$ and $\left\lbrace \mu_t\right\rbrace _{t=0}^{\infty}$ are the sequences of Lagrange multipliers associated with the sequences of constraints~(\ref{e:constraint1}) and~(\ref{e:constraint2}). For all $t\ge 0$, the first-order condition of the problem with respect to $c_{t}$ is

\[
\frac{\partial \mathcal{L}}{\partial c_{t}} = \frac{d u}{d c}(c_{t})- \lambda_{t}=0,
\]
which yields a first optimality condition:

\[
\frac{d u}{d c}(c_{t})=\lambda_{t}.
\]
For all $t\ge 0$, the first-order condition of the problem with respect to $a_{t+1}$ is

\[
\frac{\partial \mathcal{L}}{\partial a_{t+1}} = -\lambda_{t}-\mu_{t}+ \beta (1+r) \lambda_{t+1} =0,
\]
which yields a second optimality condition:

\[
\lambda_{t}+\mu_{t}=(1+r) \beta \lambda_{t+1}.
\]
In addition, the complementary slackness conditions impose that for all $t\ge 0$, $\mu_{t}\ge 0$ and

\[
\mu_{t} a_{t+1}=0.
\]

These optimality conditions and the complementary slackness conditions are necessary and sufficient if the problem is well behaved. Since $c_{t}\ge 0$, the wealth $a_{t+1}$ never falls to zero because if $a_{T}=0$, then $c_{t}=0$ for all $t\ge T$, which cannot be optimal because $u'(c_{t})=\infty$ for all $t\ge T$. Hence for all $t\ge 0$, $a_{t+1}>0$ and $\mu_{t}=0$. We infer that 

\[
\lambda_t=(1+r) \beta \lambda_{t+1}
\]
for all $t\ge 0$ such that consumption $c_{t}$ satisfies the following intertemporal condition for all $t\ge 0$:
\wcSpan*[id=equation-3]{}
\setcounter{tpequation}{2}\refstepcounter{tpequation}\label{e:eulerdp}
\[
\begin{align*}
\frac{d u}{d c}(c_{t})= \beta (1+r) \frac{d u}{d c}(c_{t+1}).\tag{3}
\end{align*}
\]
This intertemporal condition characterizes the optimal consumption path. It is called the \textit{Euler equation}. For all $t\ge 0$, wealth $a_{t}$ satisfies

\[
a_{t+1}=(1+r) a_{t}-c_{t}.
\]

\subsection{Dynamic-programming approach}

Instead of looking for two infinite sequences $\left\lbrace c_t\right\rbrace _{t=0}^{\infty}$ and $\left\lbrace a_t\right\rbrace _{t=0}^{\infty}$, dynamic programming looks for a time-invariant \textit{policy function} $h$ mapping wealth at the beginning of period $t$, $a_t$, into optimal consumption in period $t$, $c_t$, such that the sequence $\left\lbrace c_t\right\rbrace _{t=0}^{\infty}$ generated by iterating
\wcSpan*[id=equation-4]{}
\setcounter{tpequation}{3}\refstepcounter{tpequation}\label{e:iter1}
\setcounter{tpequation}{4}\refstepcounter{tpequation}\label{e:iter2}
\[
\begin{align*}
c_t&=h(a_t)\tag{4}\\
a_{t+1}&=(1+r) a_t-c_t,\tag{5}
\end{align*}
\]
starting from initial wealth $a_0$ solves the consumption-saving problem. Finding the policy function allows us to determine recursively the optimal sequence of consumption $\left\lbrace c_t\right\rbrace _{t=0}^{\infty}$.

Why do we want to find a policy function $h$ instead of an infinite sequence $\left\lbrace c_t\right\rbrace _{t=0}^{\infty}$? It is unclear that finding a function is easier than finding an infinite sequence. But it turns out that dynamic programming has three desirable properties:
\begin{enumerate}
\item Sometimes, dynamic programming allows us to find closed-form solutions for the policy function $h$.
\item Sometimes, dynamic programming allows us to characterize theoretical properties of the policy function $h$.
\item Various numerical methods are available to solve dynamic programs.
\end{enumerate}

\subsection{Value function}

To determine the policy function $h$ and solve this problem, we first need to solve for an auxiliary function that we call \textit{value function}. The value function $V(a)$ measures the optimal lifetime utility from consumption, starting with an initial wealth $a$. The value function is defined by
\wcSpan*[id=equation-6]{}
\setcounter{tpequation}{5}\refstepcounter{tpequation}\label{e:value}
\[
V(a)\equiv \max_{\left\lbrace c_t,a_{t+1}\right\rbrace _{t=0}^{\infty}} \sum_{t=0}^{\infty}\beta^t u(c_t)\tag{6}
\]
subject to for all $t\ge 0$

\[
\begin{align*}
a_{0}&= a\\
a_{t+1}&=(1+r) a_t-c_t\\
a_{t+1}&\ge 0.
\end{align*}
\]

To determine the value function, we use a theorem that says that the value function $V$ is the solution to a functional equation called the \textit{Bellman equation}:
\wcSpan*[id=equation-7]{}
\setcounter{tpequation}{6}\refstepcounter{tpequation}\label{e:recur}
\[
\begin{align*}
V(a)=\max_{c\in[0,(1+r) a]} u(c)+\beta V((1+r) a-c).\tag{7}
\end{align*}
\]

Not all optimization problems can be represented with a Bellman equation. The theorem applies only if the problem satisfies the Principle of Optimality. If a problem satisfies the Principle of Optimality, we say that it has a recursive structure. Section~\href{theory.tex#ref-sec-theory}{4} characterizes problems with a recursive structure. For instance, the consumption-saving problem has a recursive structure.

To understand where the Bellman equation comes from and what it means, we manipulate the value function. The value function can be expressed as

\[
V(a_0) = \max_{\lbrace 0\le c_t\le (1+r) a_t\rbrace _{t=0}^{\infty}} \sum_{t=0}^{\infty}\beta^t u(c_t),
\]
subject to for all $t\ge 0$

\[
a_{t+1} = (1+r) a_t-c_t.
\]
Eliminating $c_t$, we rewrite the value function as

\[
V(a_0)=\max_{\lbrace 0\le a_{t+1}\le (1+r) a_t\rbrace _{t=0}^{\infty}}\left(\sum_{t=0}^{\infty}\beta^t u((1+r) a_t-a_{t+1})\right).
\]
Separating the first term in the utility function from the rest of the sum, we obtain

\[
V(a_0)=\max_{\lbrace 0\le a_{t+1}\le (1+r) a_t\rbrace _{t=0}^{\infty}}\left(u((1+r) a_0-a_1)+\beta \sum_{t=1}^{\infty}\beta^{t-1} u((1+r) a_t-a_{t+1})\right).
\]
We re-index the terms in the sum:

\[
V(a_0) = \max_{\lbrace 0\le a_{t+1}\le (1+r) a_t\rbrace _{t=0}^{\infty}} \left\lbrace u((1+r) a_0-a_1)+\beta \sum_{t=0}^{\infty}\beta^{t} u((1+r) a_{t+1}-a_{t+2})\right\rbrace .
\]
We separate the maximization process in two stages: choose consumption in period 0 given wealth in period 0, and choose all future consumption given wealth in period 1. Thus, we rewrite the value function as

\[
V(a_0) = \max_{0\le a_{1}\le (1+r) a_0}\left(u((1+r) a_0-a_1)+\beta\max_{\lbrace 0\le a_{t+2}\le (1+r) a_{t+1}\rbrace _{t=0}^{\infty}}\left(\sum_{t=0}^{\infty}\beta^{t} u((1+r) a_{t+1}-a_{t+2})\right)\right).
\]
By definition, the second term is exactly the value function $V(a_{1})$ so we can simplify the equation to

\[
V(a_0) = \max_{0\le a_{1}\le (1+r) a_0}\left(u((1+r) a_0-a_1)+\beta V(a_1)\right).
\]
Since $c_{0}=(1+r) a_0-a_1$, we obtain

\[
V(a_0) = \max_{0\le c_{0}\le (1+r) a_0}\left(u(c_0)+\beta V((1+r) a_0-c_0)\right).
\]
This last equation is the Bellman equation. For any problem with a recursive structure, we can apply this procedure and obtain a Bellman equation.

\subsection{Policy function}

With the definition proposed in equation (\ref{e:recur}), we gave a recursive formulation to our optimization problem. Once we have determined the value function $V(a)$ for all $a$, we can easily solve for the optimal consumption level

\[
c^*=\operatorname*{argmax}_{c\in[0,(1+r) a]}u(c)+\beta V((1+r) a-c).
\]
$c^*$ is a function of initial wealth $a$. We define the policy function $h$ by

\[
h(a)\equiv c^*.
\]
The policy function provides a mapping from states to actions. It tells us how much should be consumed in the current period if initial wealth is $a$. The policy function allows us to determine the optimal path of consumption by iterating equations~(\ref{e:iter1}) and~(\ref{e:iter2}). Therefore, it allows us to solve the optimization problem defined in section~\ref{subsec:PROBLEM}. 

To solve the optimization problem defined in section~\ref{subsec:PROBLEM} using dynamic programming, we proceed in five steps:
\begin{enumerate}
\item Write down the Bellman equation.
\item Write down the first-order conditions of the optimization program.
\item Write down the Benveniste-Scheinkman equation to determine the derivative of the value function with respect to wealth, $a_{t}$.
\item Apply the Benveniste-Scheinkman equation to next period's wealth, $a_{t+1}$, and plug into the first-order conditions.
\item Derive the Euler equation, which summarizes the optimal intertemporal behavior of consumption, $c_{t}$.
\end{enumerate}

\subsection{Step 1: Bellman equation}

Thanks to the recursive structure of the consumption-saving problem, the value function $V$ satisfies the Bellman equation:
\wcSpan*[id=equation-8]{}
\setcounter{tpequation}{7}\refstepcounter{tpequation}\label{e:bellman}
\[
\begin{align*}
V(a)=\max_{c\in[0,(1+r) a]}u(c)+\beta V((1+r) a-c).\tag{8}
\end{align*}
\]

\begin{itemize}
\item Equation (\ref{e:bellman}) is a \textit{functional equation}: the unknown is the function $V$ itself. We assume for now that a solution to this equation does exist.
\item $a$ is a \textit{state variable}. It summarizes completely the information from the past that is needed to solve the forward-looking optimization problem.
\item $c$ is a \textit{control variable}. It is the variable to be chosen in the current period. It determines the value $a'$ of the state variable next period according to the \textit{transition equation}

\[
a'=(1+r) a-c.
\]
\item Notation: next period's variables are denoted using a ``prime''. So next period's consumption is $c'$ and next period's wealth is $a'$.
\end{itemize}

\subsection{Step 1': Rewrite the Bellman equation}

We rewrite the Bellman equation to get rid of the state variable, $a$, in the term $\beta V((1+r) a-c)$. To do so, we use tomorrow's state variable, $a'$, instead of today's consumption, $c$, as a control variable. We can substitute $a'$ for $c$ as a control variable because $a'$ and $c$ are directly related by $a'=(1+r) a-c$; given $a$, choosing $a'$ is equivalent to choosing $c$. 

Substituting $a'$ for $c$, the Bellman equation becomes
\wcSpan*[id=equation-9]{}
\setcounter{tpequation}{8}\refstepcounter{tpequation}\label{e:recursive}
\[
V(a)=\max_{a'\in[0,(1+r) a]} u((1+r) a-a')+\beta V(a').\tag{9}
\]
After this manipulation, the term $\beta V((1+r) a-c)$ has become $\beta V(a')$; it now only depends on the control variable. This property will be convenient to apply the envelope theorem later on.

\subsection{Step 2: First-order condition}

For now, let's assume that $V$ does exit and is differentiable. The first step is to take the first-order condition of the optimization program:

\[
\max_{a'\in[0,(1+r) a]} u((1+r) a-a')+\beta V(a').
\]
The first-order condition with respect to $a'$ is
\wcSpan*[id=equation-10]{}
\setcounter{tpequation}{9}\refstepcounter{tpequation}\label{e:FOC}
\[
\begin{align*}
0 = &(-1)\cdot \frac{\partial u}{\partial c}((1+r) a-a')+\beta \frac{d V}{d a}(a')\\
\frac{d u}{d c}(c) =& \beta \frac{d V}{d a}(a').\tag{10}
\end{align*}
\]

\subsection{Step 3: Benveniste-Scheinkman equation}

In the first-order condition~(\ref{e:FOC}), we do not know the derivative $d V/d a$ of the value function. Hence, the next step is to determine what the derivative $d V/d a$ of the value function is. To do so, we apply the envelope theorem to the Bellman equation~(\ref{e:recursive}), which yields
\wcSpan*[id=equation-11]{}
\setcounter{tpequation}{10}\refstepcounter{tpequation}\label{e:BS}
\[
\frac{d V}{d a}(a)=(1+r) \cdot \frac{d u}{d c}(c).\tag{11}
\]
This equation is the \textit{Benveniste-Scheinkman equation}. It holds for any $a$.

\subsection{Step 4: One step forward}

Equation (\ref{e:BS}) is valid for any state variable $a$. In particular, it is valid for next period's state variable:
\wcSpan*[id=equation-12]{}
\setcounter{tpequation}{11}\refstepcounter{tpequation}\label{e:BS2}
\[
\frac{d V}{d a}(a')=(1+r)\cdot \frac{d u}{d c}(c').\tag{12}
\]

\subsection{Step 5: Euler equation}

If we plug equation (\ref{e:BS2}) into equation (\ref{e:FOC}), we obtain the Euler equation in terms of current period's consumption $c$ and next period's consumption $c'$:
\wcSpan*[id=equation-13]{}
\setcounter{tpequation}{12}\refstepcounter{tpequation}\label{Euler}
\[
\frac{d u}{d c}(c)=(1+r) \beta \frac{d u}{d c}(c').\tag{13}
\]
Note that this Euler equation is exactly the same as that obtained with the Lagrangian method (see equation (\ref{e:eulerdp})). Indeed, the Lagrangian method and the dynamic programming method are equivalent.

Using the Euler equation, we could reduce the optimization problem of Section~\ref{subsec:PROBLEM} to a policy function problem: for all $a$, determine the policy function $h(a)$ such that

\[
\frac{d u}{d c}(h(a))=(1+r) \beta \frac{d u}{d c}(h((1+r) a-h(a))).
\]

\subsection{A closed-form solution}\label{guess}\wcSpan*[id=ref-guess]{}

One of the advantage of using dynamic programming is that we can sometimes find closed-form solutions for the value function and the policy function. Here, we can find a closed-form solution if we assume $u(c)=\ln\left(c\right)$. To find a closed-form solution, we use the method of undetermined coefficients. We conjecture that the value function takes the form
\wcSpan*[id=equation-14]{}
\setcounter{tpequation}{13}\refstepcounter{tpequation}\label{e:fun}
\[
V(a)=A+B \ln\left(a\right),\tag{14}
\]
where $A$ and $B$ are constants. The Bellman equation~(\ref{e:recursive}) becomes
\wcSpan*[id=equation-15]{}
\setcounter{tpequation}{14}\refstepcounter{tpequation}\label{e:AB}
\[
A+B \ln\left(a\right)=\max_{a'\in[0,(1+r) a]}\left(\ln\left((1+r) a-a'\right)+\beta \left[A+B \ln\left(a'\right)\right]\right).\tag{15}
\]
We first express the solution $a'$ of the maximization problem as a function of parameters $A$, $B$, $\beta$ and the state variable $a$. The first-order condition with respect to $a'$ yields
\wcSpan*[id=equation-16]{}
\setcounter{tpequation}{15}\refstepcounter{tpequation}\label{e:focAB}
\[
\begin{align*}
\frac{1}{(1+r) a-a'}&=\frac{\beta B}{a'}\\
a'&=(1+r) \cdot \frac{\beta B}{1+\beta B} \cdot a .\tag{16}
\end{align*}
\]
The optimal level of consumption, $c=(1+r) a-a'$, satisfies


\[
c=(1+r)\cdot \frac{1}{1+\beta B} \cdot a.
\]
We plug the expression~(\ref{e:focAB}) for the optimal $a'$ in the functional equation (\ref{e:AB}). We obtain

\[
\begin{align*}
A+B \ln\left(a\right)=&\ln\left(\frac{(1+r) a}{1+\beta B}\right)+\beta \left[A+B \ln\left(\frac{\beta B (1+r) a}{1+\beta B}\right)\right]\\
A+B \ln\left(a\right)=&\left[\beta A+ (1+\beta B) \ln\left(\frac{1+r}{1+\beta B}\right)+\beta B \ln\left(\beta B\right)\right] + (1+\beta B) \ln\left(a\right).
\end{align*}
\]
The above equation must hold for all $a$, so we must have
\wcSpan*[id=equation-17]{}
\setcounter{tpequation}{16}\refstepcounter{tpequation}\label{e:eqB}
\[
\begin{align*}
B&= 1+\beta B\\
B=&\frac{1}{1-\beta},\tag{17}
\end{align*}
\]
and

\[
A = \beta A+(1+\beta B) \ln\left(\frac{1+r}{1+\beta B}\right)+\beta B \ln\left(\beta B\right).
\]
Using the fact that $1+\beta B=\frac{1}{1-\beta}$ and $\beta B=\frac{\beta}{1-\beta}$, we infer that $A$ satisfies
\wcSpan*[id=equation-18]{}
\setcounter{tpequation}{17}\refstepcounter{tpequation}\label{e:eqA}
\[
\begin{align*}
(1-\beta) A=&(1+\beta B) \ln\left(\frac{1+r}{1+\beta B}\right)+\beta B \ln\left(\beta B\right)\\
(1-\beta) A=&\frac{1}{1-\beta} \left[\ln\left(1-\beta\right)+\ln\left(1+r\right)\right]+\frac{\beta}{1-\beta} \left[\ln\left(\beta\right)-\ln\left(1-\beta\right)\right]\\
A=&\frac{(1-\beta) \ln\left(1-\beta\right)+\beta \ln\left(\beta\right)+\ln\left(1+r\right)}{(1-\beta)^2}.\tag{18}
\end{align*}
\]
Equations (\ref{e:eqA}) and (\ref{e:eqB}) define the parameters of the value function we were solving for. With the values for the parameters $A$ and $B$, the functional form proposed in equation (\ref{e:fun}) actually solves the functional equation (\ref{e:AB}). Our guess was correct. Since the value function is unique (by theorem), we have found the value function.

Notice that we did not use the Benveniste-Scheinkman equation. The reason is that we assumed a functional form for the value function $V$, so we could compute the derivative $d V/d a$ directly, without resorting to the Benveniste-Scheinkman equation.

Using equation (\ref{e:eqB}), we can rewrite equation (\ref{e:focAB})

\[
a'=\beta (1+r) a.
\]
That is, the optimal behavior is to save a constant fraction $\beta$ of the invested wealth $(1+r) a$ and consume what is left. Let $c^{*}$ be the optimal consumption this period and $(a')^*=\beta (1+r) a $ be the optimal wealth to save for next period. The policy function is

\[
\begin{align*}
h(a)&=c^*=(1+r) a-(a')^*\\
h(a)&=(1-\beta) (1+r) a.
\end{align*}
\]
In this simple problem, we have been able to find a closed-form solution for the policy function. Unfortunately, it is usually not possible to find closed-form solutions to a dynamic program, and we must resort to numerical methods. 


\end{document}
